% ============================================================
%  ch1-3.tex — Introduction, Background, Proposed Algorithm
% ============================================================

\chapter{Introduction}

\section{Motivation}
In today's world of advanced robotics, intelligent bots and autonomous
systems are moving from research into everyday deployment. Automated vehicles ---
Autonomous Underwater Vehicles (AUV), Unmanned Aerial Vehicles (UAV) and Unmanned
Ground Vehicles (UGV) --- are increasingly used to perform tasks that are
\emph{dull, dirty or dangerous} for humans. A recurring requirement across these
applications is \textbf{complete-area coverage}: the vehicle must visit every
point of a target region, in minimum time and with minimum energy, while
respecting obstacles present in the environment.

This report applies that idea to a concrete wireless-networking problem. A UAV is
flown across a region to build a map of wireless signal strength. From that map
the system detects:
\begin{itemize}
  \item \textbf{Dead-zones} --- physical areas where a device cannot obtain a
        usable connection to any legitimate access point;
  \item \textbf{Rogue access points} --- unauthorised transmitters that imitate a
        trusted network in order to intercept traffic.
\end{itemize}

Dense wireless deployments across the 2.4, 5 and 6\,GHz bands, together with
private cellular and shared-spectrum networks, have made radio-frequency (RF)
planning harder rather than easier. Organisations increasingly run continuous
\emph{wireless assurance} programs over large campuses, warehouses and stadiums,
where surveying on foot with a laptop is slow and incomplete. A UAV that flies a
planned path and senses the spectrum from the air turns that survey into a
repeatable, automated measurement.

\section{Problem Statement and Contributions}
The work integrates path planning, passive wireless sensing and signal analysis
into a single pipeline. Its principal contributions are:
\begin{enumerate}
  \item An \textbf{obstacle-aware coverage planner} that culls nodes falling
        inside obstacle boundaries and orders the remainder with a
        Self-Organising Map \emph{seeded by a boustrophedon sweep}, giving fast,
        stable convergence to a short traversal.
  \item A \textbf{passive, privacy-preserving capture stack} that records 802.11
        management frames across the 2.4, 5 and 6\,GHz bands and tags each with a
        geographic position.
  \item A \textbf{continuous RF-heatmap dead-zone detector} that interpolates
        sampled RSSI/SNR into a surface and extracts both \emph{dead} and
        \emph{weak} contours, rather than a single binary ``no-signal'' test.
  \item A \textbf{multi-signal rogue access-point detector} that fuses an
        allow-list with evil-twin, security-capability-mismatch and hidden-node
        heuristics, remaining valid under WPA3 protected management frames and
        client MAC randomization.
  \item A \textbf{reproducible, open-source implementation} with a streaming
        analytics path, released as a documented Python package.
\end{enumerate}

\section{Host Organisation}

\subsection{Defence Research and Development Organisation (DRDO)}
The Defence Research and Development Organisation is an agency of the Government
of India charged with the military's research and development, headquartered in
New Delhi. Formed in 1958, DRDO operates a network of laboratories developing
technologies in aeronautics, armaments, electronics, land systems, life
sciences, materials, missiles and naval systems. It is India's largest research
organisation, employing thousands of scientists of the Defence Research \&
Development Service (DRDS) alongside technical and support personnel. DRDO has
delivered major systems in avionics, UAVs, radar, electronic-warfare, sonar, and
command-and-control.

\subsection{Research Centre Imarat (RCI)}
Research Centre Imarat is a premier DRDO laboratory located in Hyderabad,
responsible for the research and development of missile systems, guided weapons
and advanced avionics for the Indian armed forces. It was established by
Dr.~A.~P.~J.~Abdul Kalam in 1988. This internship was carried out within the
Directorate of Information Technology at RCI.

\section{Report Organisation}
\Cref{ch:coverage} reviews area-coverage and path-planning theory.
\Cref{ch:algo} presents the obstacle-aware coverage algorithm and its results.
\Cref{ch:sniff,ch:dead,ch:rogue} describe the packet-sniffing, dead-zone and
rogue-AP methodologies. \Cref{ch:alert,ch:analytics} cover the alert/feedback and
analytics subsystems, and \Cref{ch:impl} documents the UAV platform, the sensing
payload and the operating constraints of the system.

% ============================================================
\chapter{Area Coverage and Path Planning}
\label{ch:coverage}

\section{Optimum vs.\ Best Path Coverage}
Consider a set of checkpoints that an agent must reach. Many paths visit all of
them, but we want the one whose total cost --- time or distance --- is minimum,
so as to save energy and resources. A path that covers all intermediate
checkpoints at minimum cost is called the \textbf{optimum path}.

An incremental problem is to cover a whole \emph{area} optimally, where every
node of a region must be visited. When the map is known, the \emph{Path
Transform} algorithm (an extension of the distance transform) is considered
optimal. Zelinsky \emph{et al.}~\cite{zelinsky1993,zelinsky_phd} proposed path
coverage for a known environment, and Yun \emph{et al.}~\cite{yun2010} proposed a
genetic-algorithm approach for mobile robots.

\section{Cellular Decomposition}
Early work decomposed the free space into cells:
\begin{description}[leftmargin=1.6em,style=nextline]
  \item[Trapezoidal decomposition] splits the region into trapezoidal cells at
    obstacle \textsc{in}/\textsc{out}/\textsc{middle} events. It generates many
    cells and is not cost-effective.
  \item[Boustrophedon decomposition] (Choset) removes the \textsc{middle} event
    so that cells are merged, producing far fewer cells that are traversed in a
    back-and-forth, ``ox-like'' (\emph{boustrophedon}) motion, covering the whole
    region (\Cref{fig:ox}).
\end{description}

\begin{figure}[H]
\centering
\begin{tikzpicture}[scale=0.85]
  \draw[navy, line width=1pt] (0,0) rectangle (8,4);
  \foreach \x in {0.6,1.4,2.2,3.0,3.8,4.6,5.4,6.2,7.0,7.4}{
    \draw[accent, thick, -{Stealth[length=2mm]}] (\x,0.3) -- (\x,3.7);
  }
  \foreach \x/\y in {1.0/3.7,1.8/0.3,2.6/3.7,3.4/0.3,4.2/3.7,5.0/0.3,5.8/3.7,6.6/0.3}{
    \draw[accent, thick] (\x-0.4,\y) -- (\x+0.0,\y);
  }
\end{tikzpicture}
\caption{Area coverage by ox-like (boustrophedon) back-and-forth movement.}
\label{fig:ox}
\end{figure}

Jimenez \emph{et al.}~\cite{jimenez2007} covered a region containing polygonal
obstacles by decomposing it into obstacle-containing and obstacle-free
rectangular strips, then solving the traversal with a genetic algorithm. An
alternative \emph{spiral} method covers the area by moving along the periphery
and closing inwards (\Cref{fig:spiral}); however, it cannot guarantee a solution
and may get stuck in constricted maps with dead-end passages.

\begin{figure}[H]
\centering
\begin{tikzpicture}[scale=0.8]
  \draw[navy, line width=1pt] (0,0) rectangle (6,4);
  \draw[accent, thick, -{Stealth[length=2mm]}]
     (0.4,0.4) -- (5.6,0.4) -- (5.6,3.6) -- (0.4,3.6) -- (0.4,0.9)
     -- (5.1,0.9) -- (5.1,3.1) -- (0.9,3.1) -- (0.9,1.4)
     -- (4.6,1.4) -- (4.6,2.6) -- (1.4,2.6) -- (1.4,1.9) -- (4.1,1.9);
\end{tikzpicture}
\caption{Area coverage by spiral (periphery-inward) movement.}
\label{fig:spiral}
\end{figure}

\section{Heuristics for Optimum Paths}
Optimal path-covering algorithms are often not computationally feasible, so
\emph{heuristics} are used to obtain a near-optimal solution. They fall into
three families:
\begin{enumerate}
  \item \textbf{Construction} algorithms (e.g.\ 2-opt, 3-opt) build the path by
        adding one node at a time.
  \item \textbf{Improvement} algorithms (greedy 2-opt, greedy 3-opt) repeatedly
        exchange nodes to reduce the objective-function value (OFV).
  \item \textbf{Hybrid} algorithms (neural networks, simulated annealing, genetic
        algorithms) combine construction and improvement.
\end{enumerate}
Comparative studies~\cite{tsp_compare,youtube_tsp} show that for very large node
counts, neural-network and genetic methods produce the shortest paths, while
greedy methods are faster but longer. Crucially, \emph{all} of these classical
heuristics assume an \textbf{obstacle-free} path --- they cannot modify a route
when an obstacle lies between two connected nodes. Closing that gap is the aim of
the coverage algorithm in the next chapter.

% ============================================================
\chapter{Obstacle-Aware Coverage Algorithm}
\label{ch:algo}

\section{Sensing-Node Grid}
The target region is covered by a grid of nodes placed equidistantly. The spacing
is derived from wireless reception geometry: a receiver with reception radius
$r$ ``covers'' a circular region of radius $r$ centred at its position
(\Cref{fig:rx}). To tile the plane without overlap or gaps, adjacent nodes are
spaced by the diameter $2r$ of the reception circle.

\begin{figure}[H]
\centering
\begin{tikzpicture}[scale=0.7]
  \draw[navy, line width=1pt] (0,0) rectangle (9,5.2);
  \foreach \i in {0,...,5}{
    \foreach \j in {0,...,3}{
      \draw[accent] ({0.9+\i*1.4},{0.9+\j*1.1}) circle (0.52);
      \fill[accent] ({0.9+\i*1.4},{0.9+\j*1.1}) circle (1pt);
    }
  }
  \node[right, font=\small] at (9.1,4.0) {Receiver antenna, radius $r$};
  \draw[slate,-{Stealth[length=2mm]}] (9.0,4.0) -- ({0.9+5*1.4},{0.9+3*1.1});
\end{tikzpicture}
\caption{Area coverage by a receiver antenna ($R_x$); the reception circles tile
the plane when node spacing equals the diameter $2r$.}
\label{fig:rx}
\end{figure}

For a region of width \code{xdist} metres and a device reception range
\code{RxRange} metres, the number of nodes along axis $j$ (rows or columns) is
\begin{equation}
  n_j \;=\; \frac{\texttt{xdist}}{2 \cdot \texttt{RxRange}} , \qquad j \in \{\text{rows},\text{columns}\}.
\end{equation}
Each node is then assigned latitude/longitude and local $(x,y)$ coordinates.

\section{Obstacle-Aware Node Culling}
The system takes a map image as input and outputs a coverage path. Uniform nodes
are generated and stored in a \code{Universal} set; obstacle bounding coordinates
are obtained via \code{getObstacleCoordinates()}; and nodes lying inside an
obstacle's polygon are removed by \code{filterNodes()} using a ray-casting
point-in-polygon test. Only the surviving (\code{filter}) nodes are passed to the
path planner, so no route segment is ever planned through an obstacle.

\begin{algorithm}[H]
\caption{Obstacle-aware coverage map generation}
\KwIn{Map image $m$, device reception range $d$}
\KwOut{Coverage path over the free region}
\textbf{Universal} $\leftarrow$ \texttt{generateUniformNodes($m$, $d$)}\;
\textbf{Bound\_Coord} $\leftarrow$ \texttt{getObstacleCoordinates($m$)}\;
\textbf{filter} $\leftarrow$ \texttt{filterNodes(Universal, Bound\_Coord)} \tcp*{point-in-polygon cull}
\textbf{Map} $\leftarrow$ \texttt{generateCoveragePath(filter)}\;
\Return \textbf{Map}\;
\end{algorithm}

\section{Seeded Self-Organising Map Planner}
The traversal is produced by a Self-Organising Map (SOM): a one-dimensional ring
of neurons is fitted to the filtered nodes so that, once converged, reading the
neurons in ring order yields a short closed tour. A cold-started SOM converges
slowly and is sensitive to its learning-rate schedule. To avoid both problems the
ring is \textbf{initialised along a boustrophedon sweep} of the filtered nodes;
the SOM then only has to refine an already-reasonable tour, which reduces the
iterations needed and makes run-time predictable.

\begin{algorithm}[H]
\caption{\texttt{generateCoveragePath} --- seeded SOM path planner}
\KwIn{filtered node set \texttt{filter} ($k$ nodes), learning rate $\ell r$, decay $\alpha$}
seed $\leftarrow$ \texttt{boustrophedonOrder(filter)} \tcp*{ox-like sweep}
network $\leftarrow$ \texttt{initRingAlong(seed, $k \cdot n$)} \tcp*{$n$ neurons per node}
Initialise $\ell r$\;
\For{$i \leftarrow 1$ \KwTo maxIterations}{
  winnerNode $\leftarrow$ \texttt{selectClosest(network, random(filter))}\;
  \texttt{shiftNeighbourhood(winnerNode, random(filter))} \tcp*{Gaussian on the ring}
  $\ell r \leftarrow \ell r \cdot \alpha$\;
  \lIf{$\ell r <$ threshold}{\textbf{break}}
}
\Return neuron positions as ordered path\;
\end{algorithm}

At each iteration a random node is chosen, the closest neuron (\emph{winner}) is
selected, and the winner together with its ring neighbours is shifted toward the
node by a Gaussian function of ring distance. After each pass the learning rate
decays by a factor $\alpha\in(0,1)$; when $\ell r$ falls below a threshold the
process stops, leaving the neuron ring as the coverage path over the filtered
nodes.

\begin{highlight}[Why seed the SOM]
Seeding the neuron ring with a boustrophedon ordering turns the SOM from a
cold-start optimiser into a \emph{refiner}. Because the starting tour already
sweeps the region in rows, the network spends its iterations straightening
corners and routing cleanly around culled obstacle nodes instead of discovering
the global layout from scratch --- the same quality of tour for a fraction of the
iterations, and without hand-tuning the schedule per map. For small node counts
the result can be checked against an exact TSP solver, and against a
reinforcement-learning coverage policy for larger ones.
\end{highlight}

\section{Results}
\Cref{fig:bestpath} shows the coverage of a $9\times9$ node grid with one
obstacle. \Cref{fig:iter} plots how the number of training iterations grows with
node count, and \Cref{tab:res0,tab:res1} report iterations, run-time and path
length without and with a single obstacle.

\begin{figure}[H]
\centering
\begin{tikzpicture}[scale=0.42]
  \draw[navy, line width=0.8pt] (-0.5,-0.5) rectangle (9.5,9.5);
  \foreach \x in {0,...,9}\foreach \y in {0,...,9}{\fill[slate] (\x,\y) circle (1.6pt);}
  \fill[warn!30, draw=warn] (4,4) rectangle (6,6);
  \draw[accent, line width=1pt]
    (0,0)--(0,9)--(1,9)--(1,0)--(2,0)--(2,9)--(3,9)--(3,0)--(4,0)--(4,3)
    --(6,3)--(6,0)--(7,0)--(7,9)--(8,9)--(8,0)--(9,0)--(9,9)
    (4,7)--(4,9)--(6,9)--(6,7)--(4,7);
\end{tikzpicture}
\caption{Coverage of $9\times9$ nodes in a region with one obstacle (shaded). The
path sweeps columns and routes around the obstacle.}
\label{fig:bestpath}
\end{figure}

\begin{figure}[H]
\centering
\begin{tikzpicture}
\begin{scope}
  \def\xa{0}\def\xb{8}\def\ya{0}\def\yb{4.2}
  \draw[line, thin] (\xa,\ya) grid[xstep=1.6,ystep=1.05] (\xb,\yb);
  \draw[slate,-{Stealth[length=2mm]}] (\xa,\ya) -- (\xb+0.5,\ya) node[right,font=\small]{Nodes};
  \draw[slate,-{Stealth[length=2mm]}] (\xa,\ya) -- (\xa,\yb+0.4) node[above,font=\small]{Iterations};
  \draw[accent, line width=1.2pt, mark=*, mark options={fill=accent}]
    plot coordinates {
      (0.59,1.87)(1.85,2.70)(3.28,2.63)(4.78,4.09)(6.05,4.68)(8.0,5.22)};
  \foreach \x/\lab in {0.59/49,1.85/64,3.28/81,4.78/100,6.05/121,8.0/144}
    \node[below, font=\scriptsize] at (\x,\ya) {\lab};
\end{scope}
\end{tikzpicture}
\caption{Number of training iterations varying with the number of nodes
(single-obstacle case).}
\label{fig:iter}
\end{figure}

\begin{table}[H]
\centering
\caption{Algorithm results in a region with \textbf{0 obstacles}.}
\label{tab:res0}
\rowcolors{2}{mist}{white}
\begin{tabularx}{\textwidth}{@{}>{\raggedright\arraybackslash}X r r r@{}}
\toprule
\rowcolor{navy}
\textcolor{white}{\textbf{Node count}} & \textcolor{white}{\textbf{Iterations}} &
\textcolor{white}{\textbf{Time (s)}} & \textcolor{white}{\textbf{Path length}} \\
\midrule
49  & 19\,901 & 24.628574 & 86.981857 \\
64  & 20\,791 & 25.240143 & 115.396970 \\
81  & 20\,711 & 26.535462 & 147.811770 \\
100 & 22\,278 & 27.605774 & 184.226388 \\
121 & 22\,914 & 28.888857 & 224.640892 \\
144 & 23\,493 & 29.965775 & 269.055320 \\
\bottomrule
\end{tabularx}
\end{table}

\begin{table}[H]
\centering
\caption{Algorithm results in a region with a \textbf{single obstacle}.}
\label{tab:res1}
\rowcolors{2}{mist}{white}
\begin{tabularx}{\textwidth}{@{}>{\raggedright\arraybackslash}X r r r@{}}
\toprule
\rowcolor{navy}
\textcolor{white}{\textbf{Node count}} & \textcolor{white}{\textbf{Iterations}} &
\textcolor{white}{\textbf{Time (s)}} & \textcolor{white}{\textbf{Path length}} \\
\midrule
49  & 19\,901 & 24.286645 & 86.981857 \\
64  & 20\,791 & 25.017088 & 115.396962 \\
81  & 20\,711 & 26.104424 & 147.811770 \\
100 & 22\,278 & 27.374288 & 184.226388 \\
121 & 22\,914 & 28.161455 & 224.640892 \\
144 & 23\,493 & 29.608875 & 269.055320 \\
\bottomrule
\end{tabularx}
\end{table}

The iterations were observed for different node counts, with and without an
obstacle. Although the generated path length is similar with and without a single
obstacle, the run-time increases by small amounts in the obstacle case, as the
learning rate is depleted below the threshold more slowly. After each set of
operations the learning rate is reduced by a factor of $\alpha$; once it falls
below the threshold the operation stops, and the path corresponding to the neuron
positions is produced.
